\Conclusion

В~результате выполнения выпускной квалификационной работы были решены все поставленные задачи.

\begin{enumerate}
  \item \textbf{Проведён анализ предметной области и~существующих решений.} Исследован процесс управления цифровым присутствием субъектов МСБ, построена модель текущего процесса (AS-IS). Выполнен сравнительный анализ семи существующих решений трёх классов (CRM, SMM, ORM) методом взвешенных критериев, который показал, что ни одно из них не обеспечивает комплексного управления консистентностью данных на разнородных платформах.

  \item \textbf{Сформулированы функциональные и~нефункциональные требования к~системе.} Определены семь функциональных требований (управление бизнес-объектами, подключение каналов, контроль консистентности, оркестрация действий, контур коммуникаций, журналирование, ролевой доступ) и~восемь нефункциональных требований в~привязке к~модели качества ISO/IEC~25010:2023. Предложена концептуальная модель данных и~описано целевое состояние процесса (TO-BE).

  \item \textbf{Обоснован выбор технологий и~архитектурных решений.} Методом взвешенных критериев обоснован выбор мультиагентной архитектуры с~гибридной моделью интеграции (API~+ RPA), языка Go для серверной части (балл 2{,}80 из~3{,}00) и~фреймворка Next.js для клиентской части (балл 3{,}00 из~3{,}00). Выполнена классификация целевых платформ по типу интеграции.

  \item \textbf{Спроектирована архитектура мультиагентной системы.} Разработана микросервисная архитектура из шести компонентов (API-сервис, оркестратор, веб-приложение, три платформенных агента) и~разделяемой библиотеки, взаимодействующих через REST, SSE и~NATS. Описаны протокол Agent-to-Agent (A2A), алгоритм агентного цикла оркестратора и~алгоритм ротации JWT-токенов.

  \item \textbf{Спроектирована модель данных.} Построена модель на принципе гетерогенного хранения: реляционная СУБД PostgreSQL (5~таблиц) для структурированных данных с~ACID-гарантиями и~документоориентированная MongoDB (5~коллекций) для данных с~гибкой схемой.

  \item \textbf{Реализован прототип программного комплекса.} Общий объём исходного кода составляет порядка 26\,000~строк. Реализованы три платформенных агента: Telegram-агент (5~инструментов, Telegram Bot API), VK-агент (9~инструментов, VK API) и~Yandex.Business-агент (4~инструмента, браузерная автоматизация Playwright). Пользовательский интерфейс включает 10~экранных форм. Развёртывание контейнеризовано с~помощью Docker~Compose (9~контейнеров).

  \item \textbf{Проведено тестирование разработанного решения.} Выполнено функциональное тестирование по критическому пути пользователя для роли владельца бизнеса, подтвердившее работоспособность основных сценариев: регистрация и~аутентификация, настройка профиля бизнеса, подключение интеграций, управление данными через ИИ-ассистента, сбор отзывов и~публикация контента.
\end{enumerate}

Таким образом, цель работы~--- разработка программного комплекса на основе мультиагентной архитектуры с~гибридной моделью интеграции для управления цифровым присутствием субъектов МСБ~--- достигнута. Разработанная система позволяет управлять данными на нескольких платформах из единого интерфейса с~использованием естественного языка, при этом гибридная модель интеграции обеспечивает работу как с~платформами, предоставляющими официальные API, так и~с~платформами без публичного программного интерфейса.
